% GENERATED FILE. Edit canonical lessons, then run npm run generate:books.
% canonical-source-hash: fnv1a64:040ef928c6fa8053
% canonical-lessons: JA-C384-hikoujou, JA-C384-bin, JA-C384-ririku, JA-C384-chakuriku, JA-C384-tabidatsu

\chapter{To the Airfield}
\label{ch:ja-to-the-airfield}
\hlchaptermodality{384}

\begin{chapteropening}
\textbf{By the end of this chapter:} \emph{I can talk about flights taking off and landing.}

Complete the last lesson of chapter 384: I can talk about flights taking off and landing.
\end{chapteropening}

\section[\texorpdfstring{\ja{ひこうじょう} (hikōjō)}{ひこうじょう (hikōjō)}]{\ja{ひこうじょう} (hikōjō) — an airfield}

\label{lesson:JA-C384-hikoujou}



\begin{quote}
Before the new one: say the Japanese for swimming in the sea, then the Japanese for to be see-through.
\end{quote}

\subsection*{You'll want to know: \ja{ひこうじょう}}
\textbf{\ja{ひこうじょう}} — \emph{hikōjō} — \textquotedblleft{}an airfield\textquotedblright{}.

\begin{quote}
\textbf{\ja{この} \ja{しまには} \ja{ちいさい} \ja{ひこうじょうが} \ja{あります。}} — \emph{kono shima ni wa chiisai hikōjō ga arimasu} — \textquotedblleft{}There is a small airfield on this island.\textquotedblright{}
\end{quote}

\begin{cousinweb}
\textbf{\ja{ひこうじょう}} (\emph{hikōjō}) is Sino-Japanese (kango): three Chinese characters read with their on'yomi, readings borrowed from Chinese. They are \emph{hi}, \textquotedblleft{}fly\textquotedblright{} (Middle Chinese \emph{pj+j}), \emph{kō}, \textquotedblleft{}going\textquotedblright{} (\emph{haeng}) and \emph{jō}, \textquotedblleft{}place\textquotedblright{} (\emph{drjang}) (KANJIDIC2; Middle Chinese in Baxter's transcription).
\end{cousinweb}

\begin{grammarlens}[title={What you've built}]
One more word for getting around.
\end{grammarlens}

\subsection*{Your turn}
\begin{itemize}
\raggedright
  \item \emph{Say it:} \emph{hikōjō}
  \item \emph{Say it:} \emph{hikōjō}, once more
  \item \emph{Say it:} \emph{kono shima ni wa chiisai hikōjō ga arimasu} — There is a small airfield on this island.
  \item \emph{Recall:} say the Japanese for swimming in the sea, then the Japanese for to be see-through, then say \emph{hikōjō} again
\end{itemize}

\begin{tcolorbox}[breakable,colback=teal!4,colframe=teal!35!black,title={Before you move on}]
What does \ja{ひこうじょう} mean? (\textquotedblleft{}An airfield\textquotedblright{}.) Say it once more.
\end{tcolorbox}
\section[\texorpdfstring{\ja{びん} (bin)}{びん (bin)}]{\ja{びん} (bin) — a flight; a scheduled service}

\label{lesson:JA-C384-bin}



\begin{quote}
Before the new one: say the Japanese for to be see-through, then the Japanese for an airfield.
\end{quote}

\subsection*{You'll want to know: \ja{びん}}
\textbf{\ja{びん}} — \emph{bin} — \textquotedblleft{}a flight; a scheduled service\textquotedblright{}.

\begin{quote}
\textbf{\ja{つぎの} \ja{びんは} \ja{なんじに} \ja{でますか。}} — \emph{tsugi no bin wa nanji ni demasu ka} — \textquotedblleft{}What time does the next flight leave?\textquotedblright{}
\end{quote}

\begin{cousinweb}
\textbf{\ja{びん}} (\emph{bin}) is Sino-Japanese (kango): one Chinese character read with its on'yomi, a reading borrowed from Chinese: \emph{bin}, \textquotedblleft{}convenience\textquotedblright{} (Middle Chinese \emph{bjienH}) (KANJIDIC2; Middle Chinese in Baxter's transcription).
\end{cousinweb}

\begin{grammarlens}[title={What you've built}]
One more word for getting around.
\end{grammarlens}

\subsection*{Your turn}
\begin{itemize}
\raggedright
  \item \emph{Say it:} \emph{bin}
  \item \emph{Say it:} \emph{bin}, once more
  \item \emph{Say it:} \emph{tsugi no bin wa nanji ni demasu ka} — What time does the next flight leave?
  \item \emph{Recall:} say the Japanese for to be see-through, then the Japanese for an airfield, then say \emph{bin} again
\end{itemize}

\begin{tcolorbox}[breakable,colback=teal!4,colframe=teal!35!black,title={Before you move on}]
What does \ja{びん} mean? (\textquotedblleft{}A flight; a scheduled service\textquotedblright{}.) Say it once more.
\end{tcolorbox}
\section[\texorpdfstring{\ja{りりく} (ririku)}{りりく (ririku)}]{\ja{りりく} (ririku) — takeoff}

\label{lesson:JA-C384-ririku}



\begin{quote}
Before the new one: say the Japanese for an airfield, then the Japanese for a flight.
\end{quote}

\subsection*{You'll want to know: \ja{りりく}}
\textbf{\ja{りりく}} — \emph{ririku} — \textquotedblleft{}takeoff\textquotedblright{}.

\begin{quote}
\textbf{\ja{ひこうきは} \ja{くじに} \ja{りりくしました。}} — \emph{hikōki wa kuji ni ririku shimashita} — \textquotedblleft{}The plane took off at nine.\textquotedblright{}
\end{quote}

\begin{cousinweb}
\textbf{\ja{りりく}} (\emph{ririku}) is Sino-Japanese (kango): two Chinese characters read with their on'yomi, readings borrowed from Chinese. They are \emph{ri}, \textquotedblleft{}detach\textquotedblright{} (Middle Chinese \emph{lje}) and \emph{riku}, \textquotedblleft{}land\textquotedblright{} (\emph{ljuwk}) (KANJIDIC2; Middle Chinese in Baxter's transcription).
\end{cousinweb}

\begin{grammarlens}[title={What you've built}]
One more word for getting around.
\end{grammarlens}

\subsection*{Your turn}
\begin{itemize}
\raggedright
  \item \emph{Say it:} \emph{ririku}
  \item \emph{Say it:} \emph{ririku}, once more
  \item \emph{Say it:} \emph{hikōki wa kuji ni ririku shimashita} — The plane took off at nine.
  \item \emph{Recall:} say the Japanese for an airfield, then the Japanese for a flight, then say \emph{ririku} again
\end{itemize}

\begin{tcolorbox}[breakable,colback=teal!4,colframe=teal!35!black,title={Before you move on}]
What does \ja{りりく} mean? (\textquotedblleft{}Takeoff\textquotedblright{}.) Say it once more.
\end{tcolorbox}
\section[\texorpdfstring{\ja{ちゃくりく} (chakuriku)}{ちゃくりく (chakuriku)}]{\ja{ちゃくりく} (chakuriku) — landing (of a plane)}

\label{lesson:JA-C384-chakuriku}



\begin{quote}
Before the new one: say the Japanese for a flight, then the Japanese for takeoff.
\end{quote}

\subsection*{You'll want to know: \ja{ちゃくりく}}
\textbf{\ja{ちゃくりく}} — \emph{chakuriku} — \textquotedblleft{}landing (of a plane)\textquotedblright{}.

\begin{quote}
\textbf{\ja{ひこうきが} \ja{くうこうに} \ja{ちゃくりくしました。}} — \emph{hikōki ga kūkō ni chakuriku shimashita} — \textquotedblleft{}The plane landed at the airport.\textquotedblright{}
\end{quote}

\begin{cousinweb}
\textbf{\ja{ちゃくりく}} (\emph{chakuriku}) is Sino-Japanese (kango): two Chinese characters read with their on'yomi, readings borrowed from Chinese. They are \emph{chaku}, \textquotedblleft{}arrive\textquotedblright{} (Middle Chinese \emph{drjak}) and \emph{riku}, \textquotedblleft{}land\textquotedblright{} (\emph{ljuwk}) (KANJIDIC2; Middle Chinese in Baxter's transcription).
\end{cousinweb}

\begin{grammarlens}[title={What you've built}]
One more word for getting around.
\end{grammarlens}

\subsection*{Your turn}
\begin{itemize}
\raggedright
  \item \emph{Say it:} \emph{chakuriku}
  \item \emph{Say it:} \emph{chakuriku}, once more
  \item \emph{Say it:} \emph{hikōki ga kūkō ni chakuriku shimashita} — The plane landed at the airport.
  \item \emph{Recall:} say the Japanese for a flight, then the Japanese for takeoff, then say \emph{chakuriku} again
\end{itemize}

\begin{tcolorbox}[breakable,colback=teal!4,colframe=teal!35!black,title={Before you move on}]
What does \ja{ちゃくりく} mean? (\textquotedblleft{}Landing (of a plane)\textquotedblright{}.) Say it once more.
\end{tcolorbox}
\section[\texorpdfstring{\ja{たびだつ} (tabidatsu)}{たびだつ (tabidatsu)}]{\ja{たびだつ} (tabidatsu) — to set off on a journey}

\label{lesson:JA-C384-tabidatsu}



\begin{quote}
Before the new one: say the Japanese for takeoff, then the Japanese for landing (of a plane).
\end{quote}

\subsection*{You'll want to know: \ja{たびだつ}}
\textbf{\ja{たびだつ}} — \emph{tabidatsu} — \textquotedblleft{}to set off on a journey\textquotedblright{}.

\begin{quote}
\textbf{\ja{あには} \ja{あした} \ja{くうこうから} \ja{たびだちます。}} — \emph{ani wa ashita kūkō kara tabidachimasu} — \textquotedblleft{}My brother sets off from the airport tomorrow.\textquotedblright{}
\end{quote}

Polite \textbf{\ja{たびだちます}}, \emph{tabidachimasu}; a godan verb. Past (polite): \textbf{\ja{あにが} \ja{くうこうから} \ja{たびだちました。}}, \emph{ani ga kūkō kara tabidachimashita}.

\begin{cousinweb}
\textbf{\ja{たびだつ}} (\emph{tabidatsu}) is a native Japanese word (wago), inherited rather than borrowed from Chinese, built from \emph{tabi}, \textquotedblleft{}journey\textquotedblright{} and \emph{datsu}, \textquotedblleft{}setting out\textquotedblright{} (from \emph{tatsu}); its kanji were attached later as kun'yomi (JMdict).
\end{cousinweb}

\begin{grammarlens}[title={What you've built}]
One more word for getting around.
\end{grammarlens}

\subsection*{Your turn}
\begin{itemize}
\raggedright
  \item \emph{Say it:} \emph{tabidatsu}
  \item \emph{Say it:} \emph{tabidatsu}, once more
  \item \emph{Say it:} \emph{ani wa ashita kūkō kara tabidachimasu} — My brother sets off from the airport tomorrow.
  \item \emph{Recall:} say the Japanese for takeoff, then the Japanese for landing (of a plane), then say \emph{tabidatsu} again
\end{itemize}

\begin{tcolorbox}[breakable,colback=teal!4,colframe=teal!35!black,title={Before you move on}]
What does \ja{たびだつ} mean? (\textquotedblleft{}To set off on a journey\textquotedblright{}.) Say it once more.
\end{tcolorbox}
